\documentclass[10pt,a4paper]{amsart}
\usepackage[margin=19mm]{geometry}
\usepackage{amsmath,amssymb,mathtools}
\usepackage[scr=rsfs,cal=euler]{mathalfa}
\usepackage{xfrac}
\usepackage[unicode,hidelinks,bookmarks=false]{hyperref}
\newcommand{\ie}{i.e.\ }
\newcommand{\cf}{cf.\ }
\newcommand{\resp}{respectively\ }
\newcommand{\Subsection}{\subsection}
\providecommand{\nobreakdash}{\nobreak\hbox{-}}
\setlength{\emergencystretch}{2em}
\multlinegap0pt
\usepackage{amssymb}
\usepackage{xcolor}
\newtheorem{lem}{Lemma}
\newtheorem{proposition}{Proposition}
\newtheorem{lemma}{Lemma}
\title{Heat equations driven by mixed local-nonlocal operators with exponential nonlinearity}
\author{Dharmendra Kumar Chaurasia, Ahmad Z. Fino, Vishvesh Kumar}
\date{Source: arXiv:2605.03657v1 (CC BY 4.0)}
\begin{document}
\maketitle

\noindent\emph{Source and licence.} Heat equations driven by mixed local-nonlocal operators with exponential nonlinearity. Version arXiv:2605.03657v1 (CC BY 4.0), distributed by arXiv under the Creative Commons Attribution 4.0 International licence, http://creativecommons.org/licenses/by/4.0/, as recorded in the arXiv licence metadata for this exact version and shown on the abstract page https://arxiv.org/abs/2605.03657v1. Full licence notice: \texttt{licenses/arxiv-2605.03657-CC-BY-4.0.txt}.\par\smallskip

\noindent\textit{Editorial scope.} Counted unit: the article's lemma lemma8 (source0.tex lines 794--804). The article's complete original proof of it is delivered with it (source0.tex lines 806--840, one \texttt{proof} environment of 35 lines). No mathematics is written, repaired or re-derived: the statement and the proof are the article's own lines, in the article's own notation, and no retained line is substituted or re-flowed.  Retained uncounted as the source-local context the proof needs: lines 262--264; lines 507--527; lines 546--548; lines 555--564; lines 613--615; lines 648--651; lines 754--764; lines 789--792.  Nothing was omitted from any retained span, so the package carries no omission record.  The retained lines cite 2 works, whose entries are transcribed verbatim from the article's own bibliography source in the e-print and delivered with short editorial labels recorded in the item's reference mapping.  No retained line contains a figure environment, an image-inclusion command or drawing code, so no image is delivered and the package declares no asset dependency.  Editorial formatting only, in addition: an AMS \texttt{amsart} preamble that restates the article's own declarations and macros used by the retained lines, namely the environments \texttt{lem}, \texttt{proposition}, \texttt{lemma} on separate counters, together with the standard packages those lines need.  The retained environments print in this document as Lemma 1, Proposition 1, Lemma 1 in order of appearance, whereas the article prints the same items with the numbering of its own full document.  The complete native source of the article as distributed in the arXiv e-print for this exact version is bundled unchanged as \texttt{sources/199803/source0.tex}.
 \begin{equation}\label{eq3}
  f(0)=0,\qquad |f(u)-f(v)|\leq C|u-v|(e^{\lambda |u|^p}+e^{\lambda |v|^p}),\quad \hbox{for all}\,\, u,v\in\mathbb{R}, 
  \end{equation} 

\begin{lem}[Inclusion properties]\label{lemma1} We have the following inclusion relations:

\begin{itemize}
\item[({1})] 
\cite[Lemma 2.3]{Majdoub2}
For every $1\leq q\leq p$, we have $$L^q(\mathbb{R}^d)\cap L^\infty(\mathbb{R}^d)\hookrightarrow \exp L^p_0(\mathbb{R}^d)\hookrightarrow \exp L^p(\mathbb{R}^d).$$ More precisely, we have 
\begin{equation}\label{eq6}
\|u\|_{\exp L^p}\leq\frac{1}{(\ln 2)^{1/p}}(\|u\|_q+\|u\|_\infty).
\end{equation}
\item[({2})] \cite[Lemma 2.3]{FK20}${}$
Let $\phi(\tau)=e^{\tau^p}-1-\tau^p$, $p>1$. For every $ q\leq 2p$, we have $L^q(\mathbb{R}^d)\cap L^\infty(\mathbb{R}^d)\hookrightarrow L^\phi_0(\mathbb{R}^d)\hookrightarrow L^\phi(\mathbb{R}^d)$, more precisely
$$
\|u\|_{ L^\phi}\leq C(p)(\|u\|_q+\|u\|_\infty).
$$
\item[({3})] \cite[Lemma 2.4]{Majdoub2}
For every $1\leq p\leq q<\infty$, we have $\exp L^p(\mathbb{R}^d)\hookrightarrow L^q(\mathbb{R}^d)$, more precisely
\begin{equation}\label{eq8}
\|u\|_q\leq\left(\Gamma\left(\frac{q}{p}+1\right)\right)^{1/q} \|u\|_{\exp L^p}.
\end{equation}
\end{itemize}
\end{lem}

\begin{equation} \label{pqbound}
\Vert e^{-t \mathcal{L}} \varphi \Vert_{r} \leq C\, t_s^{-\frac{d}{2s}(\frac{1}{q}-\frac{1}{r})}\Vert \varphi \Vert_{q},\quad t>0,
\end{equation}

\begin{proposition}\label{techprop1}${}$
Let $1\leq q\leq p$, $1\leq r\leq\infty$, and $0<s< 1$. Define $t_s:=\max\{t^s,t\}$. Then the heat semigroup $e^{-t \mathcal{L}}$ satisfies the following mapping properties:
\begin{enumerate}
\item $\left\|e^{-t\mathcal{L}} \varphi\right\|_{\exp L^p }\leq \left\|\varphi\right\|_{\exp L^p }$, \,\,\, for all $t>0,$\, $\varphi\in {\exp L^p(\mathbb{R}^d)}$.

\item $\left\|e^{-t\mathcal{L}} \varphi\right\|_{\exp L^p }\lesssim t_s^{-\frac{d}{2s q}}\left(\ln(t_s^{-\frac{d}{2s}}+1)\right)^{-1/p}\left\|\varphi\right\|_{q}$, \,\,\, for all $t>0,$\, $\varphi\in L^q(\mathbb{R}^d)$.

\item $\left\|e^{-t\mathcal{L}} \varphi\right\|_{\exp L^p}\lesssim \frac{1}{(\ln 2)^{1/p}}\left[t_s^{-\frac{d}{2s r}}\left\|\varphi\right\|_{r}+\left\|\varphi\right\|_{q}\right]$, \, for all $t>0,$\, $\varphi\in L^r(\mathbb{R}^d)\cap L^q(\mathbb{R}^d)$.
\end{enumerate}
\end{proposition}

\begin{proposition}\label{prop2}
For any $\varphi\in \exp L^p_0(\mathbb{R}^d)$, the map $t\mapsto e^{-t\mathcal{L}}\varphi$ belongs to $C([0,\infty),\exp L^p_0(\mathbb{R}^d)).$
\end{proposition}

\begin{lemma}\label{lemma10}
Let $X$ be a Banach space and $ g\in L^1(0,T;X)$, then $$\displaystyle \int_0^te^{-(t-\tau)\mathcal{L}}g(\tau)\,d\tau\in C([0,T],X).$$ Moreover, we have
$$\left\|\int_0^te^{-(t-\tau)\mathcal{L}}g(\tau)\,d\tau\right\|_{L^\infty(0,T;X)}\leq \|g\|_{L^1(0,T;X)}.$$
\end{lemma}

\begin{lemma}${}$\label{lemma7} 
Let $s\in(0,1)$, $p>1$ and $v_0\in L^p(\mathbb{R}^d)\cap L^\infty(\mathbb{R}^d)$. Suppose that the nonlinearity $f$ satisfies condition \emph{(\ref{eq3})}. Then, there exist a time $T=T(v_0)>0$ and a mild solution $v\in C([0,T],\exp L^p_0(\mathbb{R}^d))\cap L^\infty(0,T;L^\infty(\mathbb{R}^d))$ of 
\begin{equation}\label{eq24}
\left\{\begin{array}{ll}
\,\, \displaystyle {\partial_t v+\mathcal{L}v =f(v),} &\displaystyle {t>0,x\in {\mathbb{R}^d},}\\
{}\\
\displaystyle{v(0)=  v_0 \qquad\qquad}&\displaystyle{x\in {\mathbb{R}^d}}.
\end{array}
\right.
\end{equation} 
 \end{lemma}

\begin{lemma}\cite[Lemma 4.4]{Majdoub2}.\label{lemma9}
Let $v\in L^\infty(0,T;L^\infty(\mathbb{R}^d))$ for some $T>0$. Suppose that the nonlinearity $f$ satisfies condition \emph{(\ref{eq3})} with some $\lambda>0.$ Let $1<p\leq q<\infty$, and $w_1,w_2\in \exp L^p(\mathbb{R}^d)$ with $\|w_1\|_{\exp L^p},\|w_2\|_{\exp L^p}\leq M$ for sufficiently small $M>0$ (namely $2^p\lambda q M^p\leq1$). Then, there exists a constant $C=C(q)>0$ such that 
$$\|f(w_1+v)-f(w_2+v)\|_q\leq Ce^{2^{p-1}\lambda\|v\|_\infty^p}\|w_1-w_2\|_{\exp L^p}.$$
 \end{lemma}

\begin{lemma}${}$\label{lemma8}
Let $s\in(0,1)$, $p>1$, and $w_0\in \exp L^p_0(\mathbb{R}^d)$. Suppose that the nonlinearity $f$ satisfies condition \emph{(\ref{eq3})}. Let $T>0$ and $v\in L^\infty(0,T;L^\infty(\mathbb{R}^d))$ be given by Lemma $\ref{lemma7}$. Then, for $\|w_0\|_{\exp L^p}\leq\varepsilon$, with $\varepsilon\ll1$ small enough, there exist a time $\widetilde{T}=\widetilde{T}(w_0,\varepsilon, v)>0$ and a mild solution $w\in C([0,\widetilde{T}],\exp L^p_0(\mathbb{R}^d))$ to the problem 
\begin{equation}\label{eq25}
\left\{\begin{array}{ll}
\,\, \displaystyle {\partial_t w+\mathcal{L}w =f(w+v)-f(v),} &\displaystyle {t>0,x\in {\mathbb{R}^d},}\\
{}\\
w(0)= w_0, & \displaystyle{ x\in {\mathbb{R}^d}.}
\end{array}
\right.
\end{equation} 
 \end{lemma}

\begin{proof} This lemma is also proved by applying the Banach fixed-point theorem. For $\widetilde{T}>0$, we define
$$W_{\widetilde{T}}:=\left\{w\in C([0,\widetilde{T}],\exp L^p_0(\mathbb{R}^d)):\,\,\|w\|_{L^\infty(0,\widetilde{T};\exp L^p_0)}\leq 2\varepsilon\right\},$$
and we introduce the mapping $\Phi$ on $ W_{\widetilde{T}}$ by
$$\Phi(w):=e^{-t\mathcal{L}}w_0+\int_0^te^{-(t-\tau)\mathcal{L}}\left(f(w(\tau)+v(\tau))-f(v(\tau))\right)\,d\tau.$$
We first show that $\Phi$ is a contraction on $W_{\widetilde{T}}$ for sufficiently small $\varepsilon$ and $\widetilde{T}>0$. Let $w_1,w_2\in W_{\widetilde{T}}$. Using Lemma \ref{lemma1}-(1), we obtain
\begin{equation}\label{eq27}
\|\Phi(w_1)-\Phi(w_2)\|_{\exp L^p}\leq \frac{1}{(\ln 2)^{1/p}}\left(\|\Phi(w_1)-\Phi(w_2)\|_{p}+\|\Phi(w_1)-\Phi(w_2)\|_{\infty}\right).
\end{equation}
Let $r>0$ be a parameter chosen such that $r>\max\{p,\frac{d}{2}\}$. Then, by the $L^r-L^\infty$ estimate (\ref{pqbound}), we have
$$\|\Phi(w_1)-\Phi(w_2)\|_{\infty}\leq C\int_0^t(t-\tau)_s^{-\frac{d}{2 s r}}\|f(w_1(\tau)+v(\tau))-f(w_2(\tau)+v(\tau))\|_r\,d\tau.$$
Applying Lemma \ref{lemma9} and assuming that $2^p\lambda r (2\varepsilon)^p\leq1$, together with the estimate $(t-\tau)_s\geq (t-s)^s$, we obtain
\begin{eqnarray}\label{eq28}
\|\Phi(w_1)-\Phi(w_2)\|_{\infty}&\leq& Ce^{2^{p-1}\lambda\|v\|_\infty^p}\left(\int_0^t(t-\tau)^{-\frac{d}{2 r}}\,d\tau\right)\|w_1-w_2\|_{L^\infty(0,\widetilde{T};\exp L^p)}\nonumber\\
&\leq& Ce^{2^{p-1}\lambda\|v\|_\infty^p}\widetilde{T}^{1-\frac{d}{2 r}}\|w_1-w_2\|_{L^\infty(0,\widetilde{T};\exp L^p)}.
\end{eqnarray}
In a similar way, under the condition $2^p\lambda p (2\varepsilon)^p\leq1$, we get
\begin{eqnarray}\label{eq29}
\|\Phi(w_1)-\Phi(w_2)\|_{p}&\leq& \int_0^t\left\|f(w_1(\tau)+v(\tau))-f(w_2(\tau)+v(\tau))\right\|_{p}\,d\tau\nonumber\\
&\leq& Ce^{2^{p-1}\lambda\|v\|_\infty^p}\widetilde{T}\|w_1-w_2\|_{L^\infty(0,\widetilde{T};\exp L^p)}.
\end{eqnarray}
Combine estimates (\ref{eq28}) and (\ref{eq29}) with (\ref{eq27}), we deduce that, for $\varepsilon\ll1$ sufficiently small,
\begin{equation}\label{eq30}
\|\Phi(w_1)-\Phi(w_2)\|_{\exp L^p}\leq \frac{1}{2}\|w_1-w_2\|_{L^\infty(0,\widetilde{T};\exp L^p)},
\end{equation}
provided that $\widetilde{T}\ll1$ is chosen sufficiently small so that $Ce^{2^{p-1}\lambda\|v\|_\infty^p}\left(\widetilde{T}+\widetilde{T}^{1-\frac{d}{2 r}}\right)\leq 1/2$.\\
We now show that $\Phi$ maps $W_{\widetilde{T}}$ into itself. Let $w\in W_{\widetilde{T}}$. Since $w_0\in L^p(\mathbb{R}^d)\cap L^\infty(\mathbb{R}^d)$, Lemma \ref{lemma1}-(1) ensures that $w_0\in\exp L^p_0(\mathbb{R}^d)$. Consequently, by Proposition \ref{prop2}, 
$$e^{-t\mathcal{L}}w_0\in C([0,T],\exp L^p_0(\mathbb{R}^d)).$$ 
Next, applying estimates (\ref{eq28})-(\ref{eq29}) with $w_1=w$ and $w_2=0$, and assuming that $2^p\lambda r (2\varepsilon)^p\leq1$, we deduce that
$$\Phi(w)-e^{-t\mathcal{L}}w_0\in L^\infty(0,T;\exp L^p_0(\mathbb{R}^d)),$$
thanks to the embedding $L^p(\mathbb{R}^d)\cap L^\infty(\mathbb{R}^d)\hookrightarrow \exp L^p_0(\mathbb{R}^d)$ (Lemma \ref{lemma1}-(1)). By the smoothing effect of the fractional semigroup $e^{-t\mathcal{L}}$ (Lemma \ref{lemma10}), it follows that 
$$\Phi(w)-e^{-t\mathcal{L}}w_0\in C([0,T],\exp L^p_0(\mathbb{R}^d)),$$
and hence $\Phi(w)\in C([0,T],\exp L^p_0(\mathbb{R}^d))$.  Moreover, using Proposition \ref{techprop1} and estimate (\ref{eq30}) with $w_1=w$ and $w_2=0$, we obtain for $T>0$ sufficiently small that 
$$\|\Phi(w)\|_{W_{\widetilde{T}}}\leq \|w_0\|_{\exp L^p} +\frac{1}{2}\|w\|_{L^\infty(0,\widetilde{T};\exp L^p)}\leq 2\varepsilon.$$
This establishes that $\Phi(w)\in W_{\widetilde{T}}$. The conclusion then follows from the Banach fixed-point theorem. This completes the proof of Lemma \ref{lemma8}.
\end{proof}

\begin{thebibliography}{99}
\bibitem[FK20]{FK20} A. Z. Fino and M. Kirane, The Cauchy problem for heat equation with fractional Laplacian and exponential nonlinearity, \emph{Commun. Pure Appl. Anal.} \textbf{19} (2020), 3625--3650.
\bibitem[Majdou]{Majdoub2} \newblock M. Majdoub and S. Tayachi, \newblock Well-posedness, Global existence and decay estimates for the heat equation with general power-exponential nonlinearities, \newblock \emph{Proc. Int. Cong. of Math., Rio de Janeiro} \textbf{2} (2018), 2379--2404.
\end{thebibliography}
\end{document}
